\section{Architectural Motivation: Reuse, Novelty, and Two-Phase Routing}
\label{sec:motivation}

Before introducing the entropy objective, we document two routing phenomena that motivated
the SR-Core design.
Both observations come from analysis of earlier recursive-routing variants in which each
step $r$ was free to select independently; this per-step freedom is precisely what the
SR-Core constraint removes.
The findings are reported here as \emph{motivation} for that constraint and for the entropy
objective, not as properties of the final SR-Core model (where $S_t$ is identical across
all steps by construction).

\subsection{Two-Phase Routing Structure}
\label{sec:two-phase}

To diagnose the routing dynamics that motivated the SR-Core constraint, we analyzed earlier
recursive-routing variants in which each step $r$ independently selected its own block set
--- the free-routing regime that SR-Core subsequently replaces with a single shared
selection.
In those variants we measure the \textbf{Jaccard similarity} between the selected block
sets of consecutive tokens at each recursion step.
Formally, for tokens at positions $t$ and $t+1$, the hard-overlap Jaccard at step $r$ is
\begin{equation}
  J_r = \frac{|S_r(t) \cap S_r(t+1)|}{|S_r(t) \cup S_r(t+1)|}
  \label{eq:jaccard}
\end{equation}
where $S_r(t) \subseteq \{1, \ldots, n\}$ is the top-$k$ selection at position $t$, step $r$.

In trained models we observe a stable two-phase pattern.
At the \textbf{entry step} ($r = 1$), Jaccard similarity between consecutive tokens is low
($J_1 \approx 0.20$): each token routes through a largely distinct block set, functioning
as a token-specific gatekeeper that selects an entry path into the block bank.
At all subsequent steps ($r = 2, \ldots, R$), Jaccard similarity collapses to a
near-constant high value ($J_{2\ldots6} \approx 0.984$): the model has settled into a
reusable core that is shared across most tokens regardless of content.

This structure emerges spontaneously from language modeling loss without any explicit routing
supervision.
It is architecturally convenient for working-set minimization --- a stable resident core
reduces cold-start misses --- but it also reveals a potential inefficiency: the reusable
core may represent over-consolidation rather than learned representational depth.

\subsection{Reuse/Novelty Tension}
\label{sec:tension}

To test whether the high-Jaccard reuse regime is optimal, we conduct
\textbf{forced-diversity ablations}: at evaluation time, we mask the $k$ blocks selected
at step $r{-}1$ before computing the top-$k$ at step $r$, forcing the model to route
through fresh blocks at each recursive step.

The results are diagnostic.
In a model trained without any routing regularization, forced diversity \emph{improves} loss
at deep recursion steps: at $r = 6$, forcing diversity reduced loss by approximately $0.104$
nats compared to the model's learned routing, which had converged to near-constant block
reuse.
This means that the high-Jaccard reuse pattern observed in trained models is not the
representationally optimal configuration --- it is a training artifact.
The routing objective (minimize next-token prediction loss) rewards early reuse because it
reduces the variance of the block computation, but this comes at the cost of the incremental
representational update that later recursive steps could provide.

These two observations together define the tension: \textbf{reuse} (same blocks, fewer
cold-start misses, cache-efficient) versus \textbf{novelty} (fresh blocks, more incremental
representation, potentially better predictions).
Router structure is causally relevant: neither unconstrained reuse nor unconstrained
diversity is a sufficient operating point.
This motivates controlled router shaping --- an objective that exposes the reuse/novelty
axis as a tunable degree of freedom rather than maximizing either endpoint.

\subsection{From Observation to Intervention}
\label{sec:intervention}

These findings motivate a router-shaping objective, but they do not prescribe the form of
that objective.
One natural candidate is a cross-token similarity penalty (e.g., encouraging lower Jaccard
between consecutive tokens' routing decisions).
We tested a soft Jaccard objective and found that it is the wrong lever: it acts on pairwise
token overlap but does not control the \emph{concentration} of the router's pre-selection
probability distribution.
In our evaluation, the soft-Jaccard variant increased route diversity (more unique block
combinations across the corpus) while simultaneously reducing hard overlap --- the opposite
of the cache-beneficial consolidation we seek.

The correct intervention point is the \textbf{sharpness of the pre-selection distribution}
at $r = 1$.
A flatter distribution means the router is uncertain which blocks to select; sharpening it
means the router assigns higher probability mass to a smaller set of blocks before the
top-$k$ cut, producing more consistent routing choices across tokens with similar
representations.
Entropy minimization acts directly on this quantity.

These observations motivate router-shaping objectives: not to maximize diversity or reuse
blindly, but to expose a controllable axis between route diversity and cache locality.
